%! TeX root = ../main.tex

\subsection{Análisis de observabilidad del sistema}
La matriz de observabilidad de un sistema se define de la forma

\begin{equation}
  \mathcal{O} = \begin{bmatrix}
    \mathbf{C} \\ \mathbf{C}\mathbf{A} \\ \vdots \\ \mathbf{C}\mathbf{A}^{n-1}
  \end{bmatrix} 
\end{equation}

Debido a que el sistema corresponde al orden 3, su matriz de observabilidad
está dada por la expresión \eqref{eq:MatrizObservabilidad}

\begin{equation}
  \mathcal{O} = \begin{bmatrix}
    \mathbf{C} \\ \mathbf{C}\mathbf{A} \\ \mathbf{C}\mathbf{A}^{2}
  \end{bmatrix} 
  \label{eq:MatrizObservabilidad}
\end{equation}

\begin{equation*}
  \begin{aligned}
    \mathbf{CA} &= 
    \begin{bmatrix}
      0 &
      0 &
      1 \\
    \end{bmatrix}
    \begin{bmatrix}
      -\frac{R}{L}  & 0 & -\frac{K_e}{L} \\
      0             & 0 &  1             \\
      \frac{K_t}{J} & 0 & -\frac{B}{J}   \\
    \end{bmatrix} \\
     &= 
    \begin{bmatrix}
      \frac{K_t}{J} & 0 & -\frac{B}{J}   \\
    \end{bmatrix}
  \end{aligned}
\end{equation*}

\begin{equation*}
  \begin{aligned}
    \mathbf{CA^2} &=
    \begin{bmatrix}
      0 & 0 & 1
    \end{bmatrix}
    \begin{bmatrix}
       \tfrac{R^2}{L^2} - \tfrac{K_e K_t}{J L}  & 0 &  \tfrac{K_e}{L}\left( \tfrac{R}{L} + \tfrac{B}{J} \right) \\[4pt]
       \tfrac{K_t}{J}                           & 0 & -\tfrac{B}{J}                                             \\[4pt]
       -\tfrac{K_t}{J}\left( \tfrac{R}{L} + \tfrac{B}{J} \right) & 0 & \tfrac{B^2}{J^2} - \tfrac{K_e K_t}{J L}
    \end{bmatrix} \\
     &=
    \begin{bmatrix}
       -\tfrac{K_t}{J}\left( \tfrac{R}{L} + \tfrac{B}{J} \right) & 0 & \tfrac{B^2}{J^2} - \tfrac{K_e K_t}{J L}
    \end{bmatrix}
  \end{aligned}
\end{equation*}

\begin{equation}
  \mathcal{O} = 
  \begin{bmatrix}
    0 & 0 & 1 \\[4pt]
    \tfrac{K_t}{J} & 0 & -\tfrac{B}{J} \\[4pt]
    -\tfrac{K_t}{J}\left( \tfrac{R}{L} + \tfrac{B}{J} \right) & 0 & \tfrac{B^2}{J^2} - \tfrac{K_e K_t}{J L}
  \end{bmatrix}
  \label{eq:MatrizObservabilidadMotor}
\end{equation}

\subsubsection{Reducción de la matriz de observabilidad}

\begin{equation*}
  \mathcal{O} = 
    \begin{bmatrix}
    0                    & 0 & 1 \\
    \tfrac{K_t}{J}       & 0 & -\tfrac{B}{J} \\
    -\tfrac{K_t}{J}\left( \tfrac{R}{L} + \tfrac{B}{J} \right) & 0 & \tfrac{B^2}{J^2} - \tfrac{K_e K_t}{JL} \\
    \end{bmatrix}
\end{equation*}

\begin{equation*}
  \begin{aligned}
    &R_1 \rightleftarrows R_3\\
    &\mathcal{O} = 
    \begin{bmatrix}
    -\tfrac{K_t}{J}\left( \tfrac{R}{L} + \tfrac{B}{J} \right) & 0 & \tfrac{B^2}{J^2} - \tfrac{K_e K_t}{JL} \\
    \tfrac{K_t}{J}       & 0 & -\tfrac{B}{J} \\
    0                    & 0 & 1 \\
    \end{bmatrix}
  \end{aligned}
\end{equation*}

\begin{equation*}
  \begin{aligned}
    &C_2 \rightleftarrows C_3\\
    &\mathcal{O} = 
    \begin{bmatrix}
    -\tfrac{K_t}{J}\left( \tfrac{R}{L} + \tfrac{B}{J} \right) & \tfrac{B^2}{J^2} - \tfrac{K_e K_t}{JL} & 0 \\
    \tfrac{K_t}{J}       & -\tfrac{B}{J} & 0 \\
    0                    & 1 & 0 \\
    \end{bmatrix}
  \end{aligned}
\end{equation*}

\begin{equation*}
  \begin{aligned}
    &R_2\leftarrow J R_2\\
    &\mathcal{O} = 
    \begin{bmatrix}
    -\tfrac{K_t}{J}\left( \tfrac{R}{L} + \tfrac{B}{J} \right) & \tfrac{B^2}{J^2} - \tfrac{K_e K_t}{JL} & 0 \\
    K_t                  & -B                     & 0 \\
    0                    &  1                     & 0 \\
    \end{bmatrix}
  \end{aligned}
\end{equation*}

\begin{equation*}
  \begin{aligned}
    &R_1\leftarrow JL R_1\\
    &\mathcal{O} = 
    \begin{bmatrix}
    -K_t\left(R + \tfrac{LB}{J}\right) & \tfrac{LB^2}{J} - K_e K_t & 0 \\
    K_t                  & -B                     & 0 \\
    0                    &  1                     & 0 \\
    \end{bmatrix}
  \end{aligned}
\end{equation*}

\begin{equation*}
  \begin{aligned}
    &R_2\leftarrow \tfrac{1}{K_t} R_2\\
    &\mathcal{O} = 
    \begin{bmatrix}
    -K_t\left(R + \tfrac{LB}{J}\right) & \tfrac{LB^2}{J} - K_e K_t & 0 \\
    1                    & -\tfrac{B}{K_t}         & 0 \\
    0                    &  1                     & 0 \\
    \end{bmatrix}
  \end{aligned}
\end{equation*}

\begin{equation*}
  \begin{aligned}
    &R_1\leftarrow \tfrac{1}{K_t} R_1\\
    &\mathcal{O} = 
    \begin{bmatrix}
      -\left(R + \tfrac{LB}{J}\right) & \tfrac{LB^2}{JK_t} - K_e  & 0 \\
    1                    & -\tfrac{B}{K_t}         & 0 \\
    0                    &  1                     & 0 \\
    \end{bmatrix}
  \end{aligned}
\end{equation*}

\begin{equation*}
  \begin{aligned}
    &R_1\leftarrow R_1 - \left(\tfrac{LB^2}{JK_t} - K_e\right) R_3\\
    &R_2\leftarrow R_2 + \tfrac{B}{K_t} R_3\\
    &\mathcal{O} = 
    \begin{bmatrix}
      -\left(R + \tfrac{LB}{J}\right) & 0                     & 0 \\
    1                    &  0                     & 0 \\
    0                    &  1                     & 0 \\
    \end{bmatrix}
  \end{aligned}
\end{equation*}

\begin{equation*}
  \begin{aligned}
    &R_1\leftarrow -\tfrac{1}{R + \tfrac{LB}{J}} R_1\\
    &\mathcal{O} = 
    \begin{bmatrix}
    1                    &  0                     & 0 \\
    1                    &  0                     & 0 \\
    0                    &  1                     & 0 \\
    \end{bmatrix}
  \end{aligned}
\end{equation*}

En este sistema, es posible observar que la matriz \eqref{eq:MatrizObservabilidadMotor}
no es de rango completo. Una explicación plausible se halla en la consideración de la
posición angular como variable de estado, la cual, al no estar relacionada, en este
modelo con las demás variables, no es observable.

\subsubsection{Matriz de observabilidad con valores experimentales}

Utilizando los mismos valores experimentales: $R = 2.1\ \Omega$, $L = 0.005\ \text{H}$, $K_t = 0.12\ \text{N·m/A}$, $K_e = 0.12\ \text{V/(rad/s)}$, $J = 0.0001\ \text{kg·m}^2$ y $B = 1\times10^{-5}\ \text{N·m·s/rad}$, la matriz de observabilidad numérica es:

\begin{equation*}
  \begin{aligned}
    &\mathcal{O}_{\text{num}} = \\
    &\begin{bmatrix}
      0 & 0 & 1 \\[2pt]
      \tfrac{0.12}{1\mathrm{e}{-4}} & 0 & -\tfrac{1\mathrm{e}{-5}}{1\mathrm{e}{-4}} \\[2pt]
      -\tfrac{0.12}{1\mathrm{e}{-4}}\left( \tfrac{2.1}{5\mathrm{e}{-3}} + \tfrac{1\mathrm{e}{-5}}{1\mathrm{e}{-4}} \right) & 0 & \tfrac{(1\mathrm{e}{-5})^2}{(1\mathrm{e}{-4})^2} - \tfrac{(0.12)^2}{1\mathrm{e}{-4} \cdot 5\mathrm{e}{-3}}
    \end{bmatrix}
  \end{aligned}
\end{equation*}

\begin{equation}
\mathcal{O}_{\text{num}} = 
\begin{bmatrix}
    0 & 0 & 1 \\[2pt]
    1.2\mathrm{e}^{3} & 0 & -1.0\mathrm{e}^{-1} \\[2pt]
    -5.04\mathrm{e}^{5} & 0 & -2.88\mathrm{e}^{4}
\end{bmatrix}
\label{eq:MatrizObservabilidadMotorValoresExperimentales}
\end{equation}
